\documentclass[11pt,a4paper]{article}
\usepackage[margin=20mm]{geometry}
\usepackage[T1]{fontenc}
\usepackage[utf8]{inputenc}
\usepackage{tikz}
\usetikzlibrary{arrows.meta,positioning}
\setlength{\parindent}{0pt}
\setlength{\parskip}{8pt}
\tikzset{step/.style={draw=blue!60!black,fill=blue!5,rounded corners,text width=10cm,align=center,minimum height=13mm,font=\small},arrow/.style={-{Latex},thick}}
\begin{document}

\begin{center}\LARGE\bfseries Increment 2: Our trainable text detector\end{center}
\textbf{Automatic runtime path}
\begin{center}\begin{tikzpicture}[node distance=8mm]
\node[step] (n0) {Prepared page image};
\node[step,below=of n0] (n1) {Letterbox and normalize pixels};
\draw[arrow] (n0) -- (n1);
\node[step,below=of n1] (n2) {Our encoder, skip connections and decoder};
\draw[arrow] (n1) -- (n2);
\node[step,below=of n2] (n3) {Text probability map};
\draw[arrow] (n2) -- (n3);
\node[step,below=of n3] (n4) {Threshold, contours, expand and remove duplicates};
\draw[arrow] (n3) -- (n4);
\node[step,below=of n4] (n5) {Perspective crops and optional own line classifier};
\draw[arrow] (n4) -- (n5);
\end{tikzpicture}\end{center}
Choose either word weights or line weights using its separate configuration.
Both use this architecture; line labels are derived by grouping FUNSD words within
entities and rows, then applying the saved page transform. Word and line weights
are trained separately. The line labels can contain fragments and are not native
line ground truth.
The detector needs our trained checkpoint. Initial random weights are not a usable OCR result. This baseline uses a fixed threshold, not a learned DB threshold branch.
All coordinates refer to pixels on the prepared image. See the README and code walkthrough for examples.
\newpage
\begin{center}\LARGE\bfseries Increment 2: development\end{center}
\begin{center}\begin{tikzpicture}[node distance=8mm]
\node[step] (n0) {Page images and labelled text polygons};
\node[step,below=of n0] (n1) {Our shrunk masks and valid-pixel mask};
\draw[arrow] (n0) -- (n1);
\node[step,below=of n1] (n2) {Our segmentation-network logits};
\draw[arrow] (n1) -- (n2);
\node[step,below=of n2] (n3) {Balanced binary cross-entropy plus Dice loss};
\draw[arrow] (n2) -- (n3);
\node[step,below=of n3] (n4) {Validate and save the best checkpoint};
\draw[arrow] (n3) -- (n4);
\end{tikzpicture}\end{center}
Development uses labelled data; runtime does not require approval of predictions. Keep train, validation and final test documents separate.
Software checks and tiny synthetic training runs do not establish real-document accuracy.
\vfill\textit{Fully Self Made OCR: architectures and algorithms implemented in this project.}
\end{document}
